\section{Simultaneous Stieltjes collisions and their finite anomalies}
\label{cl:sec}

The incoming separated-triple theorem assumes pairwise disjoint singular
grids. Its first further question asks whether simultaneous collision
constants agree with a directly merged ambient finite part
\cite[Question~1]{RepoExact}. The bilinear coincidence report supplies
such a comparison for two factors \cite{RepoCoincident}. The triple
problem contains an additional phenomenon: two singular factors can
multiply a resonant Taylor coefficient of a third regular factor.
We prove an explicit finite discrepancy, identify its dependence on the
collision rates, and then give an all-factor local completion theorem.
The distinction between cutoff and spectral prescriptions also appears
in the different endpoint setting of \cite{RepoEndpoint}; here the
collision parameter moves singularities in the real argument.

\subsection{Fixed coordinates and the geometric collision theorem}
\label{cl:geometric}

Put $f(x)=\gamma_0(x)=-\psi(x)$ and $\gamma=\gamma_0(1)$. For
$0<a<b$ and $0<\epsilon<1/b$, let
\[
 C_{a,b}(\epsilon)=\operatorname{FP}_x\int_0^1
 f(x)f(\{x+a\epsilon\})f(\{x+b\epsilon\})\,dx.
\]
At each of the three distinct grid points $0$, $1-b\epsilon$, and
$1-a\epsilon$, the finite part deletes a right interval of ambient
length $\eta$ and takes the coefficient independent of $\eta$.
The three cutoff variables can also be independent. Define the merged
moment in the same coordinate by
\begin{equation}\label{cl:merged}
 Q=\operatorname{FP}_x\int_0^1f(x)^3\,dx.
\end{equation}
This has the ordinary convergent representation
\begin{equation}\label{cl:merged-convergent}
 Q=\int_0^1\left[f(x)^3-\frac1{x^3}-\frac{3\gamma}{x^2}
       -\frac{3(\gamma^2-\zeta(2))}{x}\right]dx-\frac12-3\gamma.
\end{equation}

\begin{theorem}[Triple collision constant, with its rate dependence]
\label{cl:collision}
With $L=\log\epsilon$, define
\begin{align*}
 P_{-2}(a,b;L)&=\frac{L}{ab}
    +\frac{\log a}{a(b-a)}-\frac{\log b}{b(b-a)},\\
 P_{-1}(a,b;L)&=\gamma\left[
    \frac{L+\log a}{a}+\frac{L+\log b}{b}
    +\frac{L+\log(b-a)}{b-a}\right],\\
 P_0(a,b;L)&=\zeta(2)\left[
    \frac{a-b}{a}(L+\log a)
    +\frac{b-a}{b}(L+\log b)
    +\frac b{b-a}(L+\log(b-a))\right].
\end{align*}
Then
\begin{equation}\label{cl:rate-expansion}
 C_{a,b}(\epsilon)=\epsilon^{-2}P_{-2}(a,b;L)
    +\epsilon^{-1}P_{-1}(a,b;L)+P_0(a,b;L)+Q
    +O(\epsilon|\log\epsilon|).
\end{equation}
Consequently its coefficient independent of $\epsilon$ and $\log\epsilon$
is $Q+\zeta(2)D(a,b)$, where
\[
 D(a,b)=\frac{a-b}{a}\log a+\frac{b-a}{b}\log b
               +\frac b{b-a}\log(b-a).
\]
In particular,
\begin{align}
 C_{1,2}(\epsilon)
 ={}&\frac{L-\log2}{2\epsilon^2}
   +\frac{\gamma(5L+\log2)}{2\epsilon}
   +\frac{\zeta(2)}2(3L+\log2)+Q
   +O(\epsilon|\log\epsilon|),\label{cl:equally-spaced}\\
 \operatorname{CT}_{\epsilon}C_{1,2}(\epsilon)
 ={}&Q+\frac{\zeta(2)}2\log2\ne Q.\label{cl:finite-anomaly}
\end{align}
Here $\operatorname{CT}_{\epsilon}$ removes divergent powers and all
nonconstant logarithmic terms; it does not remove finite constants
inside a chosen collision counterterm. Removing the complete $P_0$,
including its finite part, does recover $Q$.
\end{theorem}

\begin{proof}
Choose a fixed small $\delta>0$ and work in the circle coordinate
$x\in(-\delta,\delta)$ around the merged point. The Hurwitz shift
identity \cite[Equation~25.11.3]{DLMFHurwitz} gives the exact local representation
\[
 f(\{x\})=\frac{\mathbf 1_{x>0}}x+h(x),\qquad
 h(x)=f(1+x)=\gamma-\zeta(2)x+x^2R(x),
\]
where $R$ is analytic near zero. For the three rates $(0,a,b)$ write
$S_c(x)=\mathbf 1_{x+c\epsilon>0}/(x+c\epsilon)$ and
$H_c(x)=h(x+c\epsilon)$. Expand the product of the three $S_c+H_c$.
The all-regular term and the three terms with exactly one $S_c$ are
analytic in $\epsilon$ after the prescribed one-sided finite parts
are taken. Their limits are their merged counterparts. The integral
outside this neighborhood is also analytic for small $\epsilon$.

For a term with two singular factors, let their rates be $c<d$ and
let the remaining regular factor have rate $e$. Set
$t=x+c\epsilon$, $\Delta=d-c$, $\rho=e-c$ and $\ell=\delta+c\epsilon$.
Its local integral is
\[
 \operatorname{FP}\int_0^\ell
     \frac{h(t+\rho\epsilon)}{t(t+\Delta\epsilon)}\,dt.
\]
The two necessary elementary primitives are
\begin{align*}
 A_0&=\operatorname{FP}\int_0^\ell\frac{dt}{t(t+\Delta\epsilon)}
 =\frac{\log\ell-\log(\ell+\Delta\epsilon)
                  +\log(\Delta\epsilon)}{\Delta\epsilon},\\
 A_1&=\int_0^\ell\frac{dt}{t+\Delta\epsilon}
 =\log(\ell+\Delta\epsilon)-\log(\Delta\epsilon).
\end{align*}
The constant and linear parts of $h$ contribute
$\gamma A_0-\zeta(2)(A_1+\rho\epsilon A_0)$. As $\epsilon\downarrow0$,
this equals
\[
 \frac{\gamma\log(\Delta\epsilon)}{\Delta\epsilon}
 +\zeta(2)\frac{d-e}{d-c}\log(\Delta\epsilon)
 -\frac\gamma\delta-\zeta(2)\log\delta+O(\epsilon|\log\epsilon|).
\]
The last two displayed constants are exactly the merged finite parts
of $\gamma x^{-2}-\zeta(2)x^{-1}$ on $(0,\delta)$.
The analytic quadratic remainder converges to its merged ordinary
integral. More precisely, use
\[
 \frac{(t+\rho\epsilon)^2}{t(t+\Delta\epsilon)}
 =1+\frac{(2\rho-\Delta)\epsilon}{t+\Delta\epsilon}
       +\frac{\rho^2\epsilon^2}{t(t+\Delta\epsilon)}
\]
and the analytic, locally bounded function $R$ to obtain an error
$O(\epsilon|\log\epsilon|)$. Thus the three pairs give exactly
$\epsilon^{-1}P_{-1}+P_0$ beyond their merged values.

The three-singular term is supported on $x>0$. Partial fractions give
\[
 \frac1{x(x+a\epsilon)(x+b\epsilon)}
 =\frac1{ab\epsilon^2x}
  -\frac1{a(b-a)\epsilon^2(x+a\epsilon)}
  +\frac1{b(b-a)\epsilon^2(x+b\epsilon)}.
\]
Its lower-end contribution is exactly $\epsilon^{-2}P_{-2}$. The
upper-end expression tends to $-1/(2\delta^2)$, which is precisely
$\operatorname{FP}\int_0^\delta x^{-3}\,dx$; its remaining error is
$O(\epsilon)$. Summing all eight terms and the nonsingular exterior
integral proves the formula. Only simple logarithms occur because the
zero-index singular pieces are rational functions with simple linear
factors; no $\log^2\epsilon$ term has been suppressed.
\end{proof}

The rate dependence can be isolated without an unevaluated merged
constant:
\begin{equation}\label{cl:rate-difference}
 \operatorname{CT}_{\epsilon}C_{1,3}(\epsilon)
 -\operatorname{CT}_{\epsilon}C_{1,2}(\epsilon)
 =\zeta(2)\left(\log2+\frac23\log3\right).
\end{equation}
These are positive shifts, exactly as in the definition of $C_{a,b}$.
For three digamma factors all displayed correlations change sign.

\subsection{Resonant subsets and a holomorphic completion}
\label{cl:subsets}

This section supplies an all-index explanation of the new phenomenon.
Put
\[
 G(z,x)=\zeta(1-z,x)+\frac1z
     =\sum_{m\ge0}\gamma_m(x)\frac{z^m}{m!},\qquad
 c_r(z)=(-1)^r(1-z)_r,
\]
and write
\[
 g_r(z,x)=\partial_x^rG(z,x),\qquad h_r(z)=g_r(z,1).
\]
In particular $h_0(z)=\zeta(1-z)+1/z$ and
$h_r(z)=c_r(z)\zeta(r+1-z)$ for $r\ge1$. All these functions are
holomorphic near $z=0$.

Let there be $d\ge1$ factors, positive integer frequencies $p_i$, real
shifts $a_i$, and nonnegative derivative orders $r_i$. Let
$\mathcal B$ be the finite union of their singular grids on the circle,
and define the collision set at $b\in\mathcal B$ by
\[
 C_b=\{i:p_i b+a_i\in\mathbb Z\}.
\]
For $t=x-b>0$ put
\[
 H_{b,j}(z_j,t)=
 \begin{cases}
 g_{r_j}(z_j,1+p_jt),&j\in C_b,\\
 g_{r_j}(z_j,\{p_jb+a_j\}+p_jt),&j\notin C_b.
 \end{cases}
\]
For every nonempty subset $S\subseteq C_b$ set
\[
 \Lambda_S=\sum_{i\in S}z_i,\qquad
 k_S=\sum_{i\in S}(r_i+1)-1,
\]
and define the complete resonant numerator
\[
 A_{b,S}(\boldsymbol z)=
 \left(\prod_{i\in S}c_{r_i}(z_i)p_i^{z_i-r_i-1}\right)
 [t^{k_S}]\prod_{j\notin S}H_{b,j}(z_j,t).
\]
Empty products are one. For $d\ge2$, the term with
$S=\{1,\dots,d\}$ vanishes because $k_S>0$.

\begin{theorem}[All-factor local completion]
\label{cl:allfactor}
Let $I_{\boldsymbol r}(\boldsymbol z)$ be the meromorphic continuation
of the ordinary integral
\[
 \int_0^1\prod_{i=1}^d
      g_{r_i}(z_i,\{p_i x+a_i\})\,dx,
\]
initially defined, for example, by $\Re z_i>r_i+1$ for every $i$.
Then
\begin{equation}\label{cl:completion}
 J_{\boldsymbol r}(\boldsymbol z)
 =I_{\boldsymbol r}(\boldsymbol z)
       -\sum_{b\in\mathcal B}\ \sum_{\varnothing\ne S\subseteq C_b}
              \frac{A_{b,S}(\boldsymbol z)}{\Lambda_S}
\end{equation}
is holomorphic near the origin. Its coefficient
\[
 \left(\prod_i m_i!\right)[z_1^{m_1}\cdots z_d^{m_d}]
              J_{\boldsymbol r}(\boldsymbol z)
\]
is exactly the ambient-coordinate finite part of the corresponding
product of $\gamma_{m_i}^{(r_i)}$. The theorem permits arbitrary partial
and complete collisions.
\end{theorem}

\begin{proof}
The exact Hurwitz shift gives each local factor as
\[
 g_{r_i}(z_i,\{p_ix+a_i\})
 =c_{r_i}(z_i)p_i^{z_i-r_i-1}t^{z_i-r_i-1}
       +H_{b,i}(z_i,t)\quad(i\in C_b).
\]
Expand the finite product by its nonempty subsets of singular pieces.
For $S$, expand its regular complement through degree $k_S$. The
Taylor coefficient of degree $j$ integrates meromorphically with
denominator $\Lambda_S+j-k_S$. All $j<k_S$ have denominator nonzero
near the origin and already generate their coordinate finite parts.
At $j=k_S$ the term is $A_{b,S}t^{\Lambda_S-1}$.
Its meromorphic integral on $(0,\ell)$ is
$A_{b,S}\ell^{\Lambda_S}/\Lambda_S$, while its coefficientwise
coordinate finite-part generator is
$A_{b,S}(\ell^{\Lambda_S}-1)/\Lambda_S$. Their difference is the
complete numerator $A_{b,S}/\Lambda_S$, not a residue with its
parameter dependence deleted. The Taylor remainder has exponent
$\Lambda_S$ and is integrable uniformly on a sufficiently small
spectral polydisc; all fixed spectral derivatives add only logarithmic
powers. Finite summation over grid points proves holomorphy and the
coefficient interpretation, and also constructs the local meromorphic
continuation uniquely from the initial convergent integral.
\end{proof}

For two coincident factors, a product of both singular pieces has no
regular complement of the positive degree needed for resonance. Hence
the only polar divisors through the origin are $z_i=0$. For three
coincident factors, pairs of singular pieces have one regular
complement, so $z_i+z_j=0$ contributes for the first time. This is the
structural difference behind the triple collision anomaly.

\subsection{The explicit same-point triple generator}
\label{cl:triple}

\begin{corollary}[Every coincident triple Stieltjes jet]
\label{cl:triple-completion}
At unit frequencies and a common endpoint, put $R=r_1+r_2+r_3$.
The complete generator is
\begin{align}
 J_{\boldsymbol r}(\boldsymbol z)=I_{\boldsymbol r}(\boldsymbol z)
 &-\sum_{i=1}^3\frac{c_{r_i}(z_i)}{z_i}
     \sum_{u+v=r_i}\frac{h_{r_j+u}(z_j)h_{r_k+v}(z_k)}{u!v!}\notag\\
 &-\sum_{i<j}\frac{c_{r_i}(z_i)c_{r_j}(z_j)h_{R+1}(z_k)}
          {(r_i+r_j+1)!(z_i+z_j)}.\label{cl:triple-generator}
\end{align}
In each summand $\{i,j,k\}=\{1,2,3\}$.
For $r_i=0$ this is particularly short:
\begin{equation}\label{cl:triple-zero}
 J_{\boldsymbol0}=I_{\boldsymbol0}
 -\sum_i\frac{h_0(z_j)h_0(z_k)}{z_i}
 -\sum_{i<j}\frac{h_1(z_k)}{z_i+z_j}.
\end{equation}
\end{corollary}
\begin{proof}
For a singleton $S=\{i\}$, its resonant Taylor degree is $r_i$,
and Leibniz's formula gives the inner sum in the first line.
For a pair $S=\{i,j\}$, the resonant degree is $r_i+r_j+1$.
The remaining factor has derivative order $r_k$, so its Taylor
coefficient is $h_{R+1}(z_k)/(r_i+r_j+1)!$. The three-singular
subset has no regular complement and makes no resonant contribution.
Theorem~\ref{cl:allfactor} now gives both formulas.
\end{proof}
The last sum is indispensable: $h_1(0)=-\zeta(2)\ne0$.

\subsubsection{A completely specified Tornheim realization}
\label{cl:tornheim}

The classical Tornheim function is initially
\[
 T(A,B,C)=\sum_{m,n\ge1}\frac1{m^A n^B(m+n)^C}.
\]
Let $\alpha_i=z_i-r_i$ and define
\begin{align*}
 K_3(\boldsymbol\alpha)
 &=\frac{2\prod_i\Gamma(\alpha_i)}{(2\pi)^{\sum_i\alpha_i}}
   \sum_{k=1}^3
   \cos\!\left(\frac\pi2(\alpha_i+\alpha_j-\alpha_k)\right)
        T(\alpha_i,\alpha_j,\alpha_k),\\
 K_2(A,B)&=\frac{2\Gamma(A)\Gamma(B)}{(2\pi)^{A+B}}
          \cos\!\left(\frac\pi2(A-B)\right)\zeta(A+B).
\end{align*}
The classical Fourier calculation \cite{EspinosaMoll} gives these as the integrals of,
respectively, three or two factors $\zeta(1-\alpha,x)$ in an initial
absolute-convergence domain. Its six sign cones explain the three
cosines. For the regularized Stieltjes generators one has the exact
meromorphic identity
\[
 I_{\boldsymbol r}(\boldsymbol z)
 =\left(\prod_i c_{r_i}(z_i)\right)K_3(\boldsymbol\alpha)
  +\sum_{k:r_k=0}\frac{c_{r_i}(z_i)c_{r_j}(z_j)}{z_k}
                  K_2(\alpha_i,\alpha_j)
  +\frac{\mathbf 1_{\boldsymbol r=\boldsymbol0}}{z_1z_2z_3}.
\]
Terms with two scalar $1/z_i$ factors vanish because the remaining
Hurwitz factor has mean zero in the initial domain. This combined
formula, followed by the complete local subtraction above, is a
finite all-index coefficient rule. No arithmetic reduction of the
Tornheim jets to single-zeta coordinates is asserted.

Meromorphic continuation of $T$ near the required nonpositive integer
orders can be proved directly from
\[
 T(A,B,C)=\frac1{\Gamma(C)}\int_0^\infty
 t^{C-1}\operatorname{Li}_A(e^{-t})\operatorname{Li}_B(e^{-t})\,dt.
\]
On a small interval insert the classical local expansion \cite[Equation~25.12.12]{DLMFPolylog}
\[
 \operatorname{Li}_A(e^{-t})
 =\Gamma(1-A)t^{A-1}
       +\sum_{j\ge0}\zeta(A-j)\frac{(-t)^j}{j!},
\]
which is uniform on compact parameter neighborhoods avoiding positive
integer $A$. Finite subtraction and termwise integration give explicit
linear polar divisors; the remaining Mellin integral is holomorphic.
The same argument handles the colored kernels with one or both twists
equal to one. Thus an entire-order claim valid for two nontrivial
twists must be weakened to meromorphic continuation when collisions
force trivial twists. The incoming separated-grid theorem is valid
on its stated hypotheses; it is its unrestricted extrapolation that
would be incorrect.

\subsection{Spectral rays and Laurent expansion chambers}
\label{cl:spectral}

\begin{proposition}[Exact cubic spectral-ray correction]
\label{cl:rays}
Take $r_i=0$ and $z_i=c_i t$, where $c_i\ne0$ and $c_i+c_j\ne0$.
Since the completed $J$ is holomorphic, its constant is always $Q$.
The uncompleted meromorphic integral instead has constant
\begin{align}
 \operatorname{CT}_{t}I_{\boldsymbol0}(c_1t,c_2t,c_3t)
 ={}&Q+\gamma\gamma_1\sum_i\frac{c_j+c_k}{c_i}\notag\\
       &+(\zeta(2)+\zeta'(2))\sum_{i<j}\frac{c_k}{c_i+c_j}.
       \label{cl:ray-constant}
\end{align}
Here $\gamma_1$ is the ordinary first Stieltjes constant and $\zeta'$
is the spectral derivative. On the diagonal the discrepancy is
\[
 6\gamma\gamma_1+\frac32\bigl(\zeta(2)+\zeta'(2)\bigr).
\]
\end{proposition}
\begin{proof}
Use \eqref{cl:triple-zero} and the expansions
$h_0(z)=\gamma+\gamma_1z+O(z^2)$ and
$h_1(z)=-\zeta(2)+(\zeta(2)+\zeta'(2))z+O(z^2)$.
After division by $c_it$ or $(c_i+c_j)t$, precisely the linear
numerator coefficient contributes to the constant. The holomorphic
remainder contributes $J_{\boldsymbol0}(0,0,0)=Q$. For the diagonal,
set $c_1=c_2=c_3=1$.
\end{proof}
This operation concerns spectral rays and is distinct from the
geometric collision $\epsilon\to0$ proved earlier. Mixed poles also
mean that a rectangular Laurent expansion is no longer a canonical
operation unless an expansion chamber or an iterated prescription is
specified. The holomorphic completion avoids either choice.

\begin{proposition}[Three chamber constants for one derivative moment]
\label{cl:chambers}
For $\boldsymbol r=(1,0,0)$ one has
\begin{equation}\label{cl:derivative-moment}
 J_{1,0,0}(0,0,0)
 =\operatorname{FP}\int_0^1\gamma_0'(x)\gamma_0(x)^2\,dx
 =2\gamma\zeta(2)-\zeta(3).
\end{equation}
In an iterated Laurent
chamber ordering the three moduli, the raw constant is instead
\begin{equation}\label{cl:chamber-values}
 \operatorname{CT}_{\rm chamber}I_{1,0,0}
 =\begin{cases}
 \zeta(3),& |z_1|\text{ is largest},\\
 0,& |z_1|\text{ is intermediate},\\
 -\zeta(3),& |z_1|\text{ is smallest}.
 \end{cases}
\end{equation}
On the spectral diagonal the same raw integral is identically zero:
$I_{1,0,0}(t,t,t)=0$ as a meromorphic function.
\end{proposition}
\begin{proof}
The local expansion is
\[
 \gamma_0(x)=x^{-1}+\gamma-\zeta(2)x+\zeta(3)x^2+O(x^3).
\]
Consequently the coefficient of $x^0$ in $\gamma_0(x)^3$ is
$\gamma^3-6\gamma\zeta(2)+3\zeta(3)$. Integrate the exact derivative
$(\gamma_0^3)'/3$ from the cutoff to $1$ and retain its constant.
This proves \eqref{cl:derivative-moment} directly, without a spectral
regularization.

For clarity, write $I_{1,0,0}-J_{1,0,0}=S+P$, separating singleton
and pair terms in \eqref{cl:triple-generator}. Since $c_1(z_1)=z_1-1$,
these are
\begin{align*}
 S={}&\frac{z_1-1}{z_1}
       \bigl[h_1(z_2)h_0(z_3)+h_0(z_2)h_1(z_3)\bigr]\\
   &+\frac{h_1(z_1)h_0(z_3)}{z_2}
     +\frac{h_1(z_1)h_0(z_2)}{z_3},\\
 P={}&\frac{(z_1-1)h_2(z_3)}{2(z_1+z_2)}
     +\frac{(z_1-1)h_2(z_2)}{2(z_1+z_3)}
     +\frac{h_2(z_1)}{z_2+z_3}.
\end{align*}
The singleton constant is $-2\gamma\zeta(2)$ in every chamber.
Using $h_2(0)=2\zeta(3)$, the only pair terms contributing to a
constant are
\[
 \zeta(3)\frac{z_1}{z_1+z_2}
       +\zeta(3)\frac{z_1}{z_1+z_3}.
\]
In a chamber with $|z_j|<|z_1|$, expand
$z_1/(z_1+z_j)=\sum_{k\ge0}(-z_j/z_1)^k$; its constant is one.
In the reversed chamber its expansion starts with $z_1/z_j$ and
has constant zero. Thus $P$ contributes $N\zeta(3)$, where
$N$ is the number of other variables dominated by $z_1$.
Adding the holomorphic constant in \eqref{cl:derivative-moment} gives
$(N-1)\zeta(3)$, proving all three cases.

Finally, for $\Re t>2$, the function $G(t,x)$ has equal endpoint
values and a continuous first derivative on the period. Therefore
\[
 I_{1,0,0}(t,t,t)=\int_0^1
       \partial_xG(t,x)G(t,x)^2\,dx
       =\frac{G(t,1)^3-G(t,0)^3}{3}=0.
\]
Meromorphic continuation preserves this identity. In particular,
its spectral diagonal constant is zero even though its coordinate
finite part is \eqref{cl:derivative-moment}.
\end{proof}

In contrast, for $\boldsymbol r=\boldsymbol0$ and unit frequencies,
the polar numerators do not depend on the variables in their own
denominator. Their nonnegative rectangular coefficients vanish in
every standard chamber. Thus the general mixed-pole warning should
not be misread as asserting chamber dependence of the raw constant
in that special zero-derivative case; its spectral-ray constants still
differ, as shown above.

\subsection{An exact cancellation at every number of factors}

The incoming question on odd derivative order identifies equal Stieltjes
indices and one total argument derivative as its first target
\cite[Question~3]{RepoExact}. That case has a complete local answer.
The mechanism is the elementary fundamental theorem of calculus, applied
before taking the cutoff constant; keeping its endpoint term gives the
following more general identity.

\begin{theorem}[The all-factor total-derivative projection]
\label{cl:total-derivative}
Let $d\ge1$, $m_1,\ldots,m_d\ge0$, and put
$Z_0=\{i:m_i=0\}$ and $g_i(t)=\gamma_{m_i}(1+t)$. Then
\begin{align}
 &\FP\int_0^1\sum_{i=1}^d\gamma_{m_i}'(x)
                   \prod_{j\ne i}\gamma_{m_j}(x)\,dx\notag\\
 &\qquad=-\sum_{\varnothing\ne S\subseteq Z_0}
              [t^{|S|}]\prod_{j\notin S}g_j(t).
 \label{cl:total-projection}
\end{align}
The sum on the right is empty when all $m_i>0$. Each surviving
coefficient is a finite expression in ordinary Stieltjes constants and
argument derivatives of Stieltjes functions at $1$.
\end{theorem}
\begin{proof}
For every fixed cutoff, the integrand is the ordinary derivative of
$\prod_i\gamma_{m_i}(x)$. Hence its finite-part integral is
\[
 \prod_i\gamma_{m_i}(1)
       -\CT_{x\downarrow0}\prod_i\gamma_{m_i}(x).
\]
Use the exact shift
$\gamma_{m_i}(x)=x^{-1}\log^{m_i}x+g_i(x)$ and expand by the
subset $S$ of singular pieces. The empty subset contributes exactly
$\prod_i\gamma_{m_i}(1)$, canceling the upper endpoint.
A nonempty subset contributes
$x^{-|S|}\log^{\sum_{i\in S}m_i}x\prod_{j\notin S}g_j(x)$.
It has a constant in both $x$ and $\log x$ precisely from the
Taylor degree $|S|$ and only when $S\subseteq Z_0$.
This proves \eqref{cl:total-projection}. All individual finite parts
exist by the same finite local power--log expansion, so taking the
constant after the finite sum introduces no interchange issue.
\end{proof}

For equal positive indices the theorem gives the requested cancellation
at every number of factors:
\begin{equation}\label{cl:equal-index-cancellation}
 \boxed{\quad
 \FP\int_0^1\gamma_m'(x)\gamma_m(x)^{d-1}\,dx=0,
       \qquad d\ge1,\quad m\ge1.
 \quad}
\end{equation}
For the zero index, with $h(t)=\gamma_0(1+t)$ as above, the complete
remaining local value is
\begin{equation}\label{cl:zero-index-projection}
 \FP\int_0^1\gamma_0'(x)\gamma_0(x)^{d-1}\,dx
 =-\frac1d\sum_{k=1}^d\binom dk[t^k]h(t)^{d-k}.
\end{equation}
This is a finite zeta polynomial because
$h(t)=\gamma+\sum_{j\ge1}(-1)^j\zeta(j+1)t^j$ near zero.
For $d=3$ it recovers \eqref{cl:derivative-moment}; for $d=4$ it gives
\[
 \FP\int_0^1\gamma_0'(x)\gamma_0(x)^3\,dx
 =3\gamma^2\zeta(2)-3\gamma\zeta(3)-\frac{11}{4}\zeta(4).
\]
Thus the symmetric projection removes the multiple spectral terms
entirely in this one-derivative case, leaving its explicitly specified
local endpoint contribution. Higher odd derivative projections remain
a separate question.

\subsection{Repository scope and independent numerical diagnostics}
\label{cl:numerics}

The inspected snapshot is
\texttt{dcd95baeb7c0e914b138bddef6829c5b1d52b726}.
The incoming \src{Exact_Identities} report, section
\texttt{triple\_dilation.tex}, assumes pairwise disjoint grids throughout.
Its \texttt{audit\_research.tex}, Question~1, asks to remove that
restriction and compare simultaneous collision constants with merged
finite parts. The subset theorem completes its fixed-collision local
subtraction question, and the explicit $(0,\epsilon,2\epsilon)$
calculation shows that its suggested matching is not valid for plain
geometric constant extraction. A finite rate-dependent correction
is necessary. The preceding bilinear coincidence results are not
contradicted.

The supplied script \src{code/verify_collisions.py} computes the
separated moments by three independent convergent interval integrals.
For an interval of length $\ell$ starting at the singular point of one
factor, with the other two arguments $\eta_1+t,\eta_2+t$, it uses
\[
 \int_0^\ell\left[
 \frac{f(\eta_1+t)f(\eta_2+t)-f(\eta_1)f(\eta_2)}{t}
 +h(t)f(\eta_1+t)f(\eta_2+t)\right]dt
 +f(\eta_1)f(\eta_2)\log\ell.
\]
This is exactly the ambient finite part with the simple pole visibly
removed. It does not implement the collision expansion or the
Tornheim representation. The merged moment is evaluated by a local
digamma Laurent series plus ordinary quadrature, with two independent
split points. At 70-digit working precision the computed values are
\[
 Q=-1.9662627343915343406643753384116314238258689764917\ldots,
\]
and
\[
 \frac12\zeta(2)\log2=0.5700907053214263787372899294961746726083149206823\ldots.
\]
Numerical residuals tend to zero at the proved
$O(\epsilon|\log\epsilon|)$ rate for $(a,b)=(1,2),(1,3),(2,5)$.
For the equally spaced example, the residual means
$C_{1,2}(\epsilon)-\epsilon^{-2}P_{-2}
 -\epsilon^{-1}P_{-1}-P_0-Q$; the complete finite term in $P_0$
is included in this diagnostic:
\begin{center}
\begin{tabular}{@{}rr@{}}
\toprule
$\epsilon$ & Numerical residual\\
\midrule
$10^{-6}$ & $-9.332630963\times10^{-5}$\\
$10^{-8}$ & $-1.237723998\times10^{-6}$\\
$10^{-10}$ & $-1.542186195\times10^{-8}$\\
\bottomrule
\end{tabular}
\end{center}
Omitting the finite term would instead leave a limit
$\zeta(2)\log2/2$. The full record contains seventeen collision
diagnostics, the two-split merged calculation, working precision,
and the finite correction. It is reproduced from the package root by
\begin{verbatim}
python code/verify_collisions.py --output results/collisions.json
\end{verbatim}
These are floating-point diagnostics, not certified interval enclosures;
the proof is the exact local decomposition above.

The classical inputs are the Hurwitz shift, argument-derivative and
Fourier formulas \cite{DLMFHurwitz}, the local polylogarithm expansion
\cite{DLMFPolylog}, and the Hurwitz--Tornheim product relation
\cite{EspinosaMoll}. The contributions here are the complete
coordinate-normalized subtraction, the finite geometric collision
anomaly, and their exact comparison with the specified spectral
prescriptions. The statements do not claim a reduction of every
Tornheim derivative to ordinary zeta values.
